\begin{abstract}
A recent survey of zero-knowledge verifiable machine learning covers
supervised training only; no reinforcement-learning system appears. We present
a layered artifact system for zero-knowledge proof-backed verification of
offline Deep Q-Network (DQN) training over provenance-bound committed replay
data. Datasets are committed with SHA-256 Merkle trees after a replay audit.
Each relation is stated in fixed-point integer arithmetic and implemented
twice, as a Python oracle and an SP1 zkVM guest, so disagreements surface
before proving. The fragment relation derives its sampler seed from the
dataset commitment rather than accepting one, publishes a bound on admitted
values, and clips gradients component-wise; the first two close freedoms
specific to value learning, whose labels come from the model under training and
whose sampling touches little of the buffer.

On CartPole, LunarLander and PointMaze, twenty-five of twenty-six proved
configurations verify in $0.054$--$0.130$\,s with proofs of
$1.27$--$2.85$\,MB; the exception verifies Groth16 child proofs in-guest at
$68$\,s and $1.47$\,GB. Recursive aggregation places an entire $4{,}992$-step
run under one root proof bound to the dataset behind a reported policy. A
tamper benchmark of $236$ adversarial cases across $19$ categories yields no
unexpected acceptance. We report what provability costs: on
\texttt{lunarlander-random} it reaches $-152.7$ against $-215.3$ for the tuned
baseline, and an ablation isolates the target-sync interval as decisive. Ten
security statements map to relation code, backend, tests and benchmark rows.
The artifact does not claim complete DQN training soundness or Adam proofs, and
recursive aggregation requires a CUDA prover.
\end{abstract}

\begin{IEEEkeywords}
Zero-knowledge proofs, offline reinforcement learning, verifiable machine
learning, proof-of-training, Deep Q-Network, SP1 zkVM, tamper detection.
\end{IEEEkeywords}
